\documentclass[11pt]{article}
\usepackage{amsmath, amssymb}
\usepackage{geometry}
\usepackage{array}
\usepackage{booktabs}
\geometry{margin=0.7in}
\setlength{\parindent}{0pt}
\setlength{\parskip}{0.35em}
\renewcommand{\arraystretch}{1.15}

\newcommand{\cheatsheettitle}[2]{%
\begin{center}
{\LARGE #1}\\[0.4em]
{\large #2}
\end{center}
}


\begin{document}
\cheatsheettitle{Algebra I Linear Equations Cheatsheet}{Module 2}

\section*{Equation Goal}
Use inverse operations to isolate the variable while keeping both sides equal.
\[
x+a=b \Rightarrow x=b-a
\]
\[
ax=b \Rightarrow x=\frac{b}{a}\qquad (a\ne0)
\]

\section*{Multi-Step Template}
\[
a(x+b)+c=dx+e
\]
\begin{enumerate}
  \item Distribute.
  \item Combine like terms on each side.
  \item Move variable terms to one side.
  \item Move constants to the other side.
  \item Divide by the coefficient.
\end{enumerate}

\section*{Variables on Both Sides}
\[
3x+8=7x-12
\]
Move variable terms together and constants together:
\[
20=4x,\qquad x=5
\]

\section*{Special Cases}
\[
5=5 \Rightarrow \text{infinitely many solutions}
\]
\[
5=9 \Rightarrow \text{no solution}
\]
\[
x=4 \Rightarrow \text{one solution}
\]

\section*{Literal Equations}
Solve for one variable in terms of the others.
\[
A=lw \Rightarrow l=\frac{A}{w}\qquad (w\ne0)
\]
\[
y=mx+b \Rightarrow x=\frac{y-b}{m}\qquad (m\ne0)
\]

\section*{Clearing Fractions}
Multiply every term on both sides by the LCD.
\[
\frac{x}{3}+\frac{x}{4}=7
\]
\[
12\left(\frac{x}{3}\right)+12\left(\frac{x}{4}\right)=12(7)
\]

\section*{Absolute Value Equations}
\[
|A|=c
\]
If
\[
c>0,\quad A=c\text{ or }A=-c
\]
If
\[
c=0,\quad A=0
\]
If
\[
c<0,\quad \text{no solution}
\]

\section*{Quick Checks}
\begin{itemize}
  \item Distribute before moving terms.
  \item Every operation must be done to both sides.
  \item Check answers by substitution.
  \item For absolute value, split into two equations only when the outside value is nonnegative.
\end{itemize}

\end{document}
